\section{Illusion and Spiritual Warfare}

\begin{quotex}
The \emph{Jagat} the total experience of the individual in the three states of existence appears to be true so long as Brahman, the substratum, the basis of all this creation, is not realized. It is like the illusion of silver in the mother-of-pearl.
\flright{\textsc{Shankara}, \emph{Atma Bodha}}
\end{quotex}


\paragraph{We know ideas}

As we have seen, Shankara shows that our waking life is the product of a web of illusions. So without the realization of Brahman, there is no knowledge of reality. Hence what we know are merely ideas and images manufactured in the mind.

The knowledge of Brahman, in this sense, is the knowledge of the ideas in the Divine Mind, which is the substratum onto which the illusions are projected. When we know the ideas, we know the true essences of things.

\paragraph{World as prior and independent}


The most fundamental idea we believe is that ``the world exists prior to and independent of my consciousness of it.'' A limited and finite self experiences its own privations as an obstacle to its manifestation. The struggle to manifest one's possibilities is what gives the world its moral dimension.

But what was the world like before any conscious beings experienced it? As astronomers gaze into the heavens, they identify stars, planets, asteroids, comets, moons, galaxies, black holes, and so on. But before the astronomer came into being to mark out the heavens in such a way, such things did not exist as such. The universe could only be a random collection of matter and energy. A ``star'' was only a particular agglomeration of matter and energy and had no independent existence as such.

\paragraph{Science as creative act}


In many circles, science is considered as the only, or most reliable, source of knowledge. Through mathematical models, it is believed an adequate picture of the world is achieved. However, maths is only on the dividing line; on the one hand, they are not physical, yet they are not essences.

So a scientific theory is a creation of the human mind that has the property that various phenomena can be deduced from the theory. This is called ``saving the appearances.'' For example, we have the Newtonian model of the universe that for quite a while served very well in predicting future events. For a couple of centuries, people assumed it provided a true picture of the universe.

However, it was eventually overturned by the Einsteinian model, which predicted events just as well as the Newtonian, but, in additional, correctly anticipated other phenomena that could not be explained by the earlier theory. Now many people, including professional philosophers, have told me that the Newtonian model is a good ``approximation'' to the Einsteinian, so in a sense it must be ``true''. This is to misunderstand the nature of a scientific theory which, by definition, saves the appearances. Would anyone be surprised to learn that a Shakespearian and an Italian sonnet both have 14 lines? Of course not, since that is its very definition.

Furthermore, that only makes sense if you restrict maths to arithmetic. But maths also includes geometry. Geometrically, the Newtonian and Einsteinian models are wildly divergent.

\paragraph{We deduce a world}


So this is a clue. The mind creates an illusory model of the world and from that, various things are deduced from it. That is, a human world is created based on the model and its necessary deductions. Various models are in conflict with each other, with no way to adjudicate them. Yet they are all illusory and only knowledge of Brahman is the ultimate resolution.

\paragraph{Debate of the Geometers}


For example, suppose a Euclidean and a non-Euclidean geometer engage in a debate? They both have consistent worldviews, within their own assumptions. Yet, there is no way for them to decide on the issue of the fifth postulate.

\paragraph{Will to power as fact}


If the Will to Power is considered as some metaphysical theory, it is clearly false. However, Nietzsche regarded it not as a theory, but as a fact. As a fact, it is simply obvious. In our example, the Euclidean geometer could point to his successes in surveying, perspective in art, and many other applications as a demonstration of the power of his theory. The non-Euclidean could point to his successes in GPS systems.

\paragraph{Spiritual Warfare}


Spiritual warfare, then, is the battle of ideas. There are nine levels corresponding to the angelic hierarchies, in order to reach the pure essences. However, in addition to the creative force, there are destructive forces at each level. Since ideas in themselves have no power, the real battle takes place in consciousness. Until then, the idea is a possibility of manifestation, but without the human element, it does not manifest. It is therefore legitimate to see the battle on three levels. Every idea, or zeitgeist, etc., can be evaluated in these ways:

\begin{enumerate}


\item What is the prevailing idea, ideology, etc.? How does it comport with Tradition or revelation? Is it creative or destructive?

\item Who are the conscious agents of the ideology? That is, who, or what group, is aware of the ideology and actively promotes it. They know the purpose, goals, and ultimate consequences of the ideology.

\item Who are the unconscious agents of the ideology? These people or groups latch onto the ideology without understanding its source, purpose, or consequences. They have an emotional attachment but only a vague intellectual understanding. They are oblivious to any unintended consequences.
\end{enumerate}


Arguments, logic, evidence, rhetoric, and persuasion serve only as tools in spiritual warfare and can never bring ultimate victory. Instead, what is necessary is awareness and the ability to overcome illusion. Of course, one must be clear about what constitutes ``victory''.


\flrightit{Posted on 2014-05-27 by Cologero}

\begin{center}* * *\end{center}

\begin{footnotesize}
\begin{sffamily}
\texttt{Caeliger on 2017-05-03 at 19:29 said:}

Thank you for this, brother. I have just recently, after a phase of dualistic theism, perceived the faintest glimpse of the Universal. Guénon's The Multiple States of the Being and Evola's The Hermetic Tradition have helped me a lot, and I now see invisible correspondences where there was only darkness before. Forgive my enthusiasm, but when man first hears the cry of the Hidden Treasure, he is awe-struck.

\hfill

\end{sffamily}
\end{footnotesize}

